\documentclass{article}
\usepackage{mathtools} 
\usepackage{fontspec}
\usepackage[UTF8]{ctex}
\usepackage{amsthm}
\usepackage{mdframed}
\usepackage{xcolor}
\usepackage{amssymb}
\usepackage{amsmath}


% 定义新的带灰色背景的说明环境 zremark
\newmdtheoremenv[
  backgroundcolor=gray!10,
  % 边框与背景一致，边框线会消失
  linecolor=gray!10
]{zremark}{说明}

% 通用矩阵命令: \flexmatrix{矩阵名}{元素符号}{行数}{列数}
\newcommand{\flexmatrix}[4]{
  \[
  #1 = \begin{pmatrix}
    #2_{11}     & #2_{12}     & \cdots & #2_{1#4}   \\
    #2_{21}     & #2_{22}     & \cdots & #2_{2#4}   \\
    \vdots      & \vdots      & \ddots & \vdots     \\
    #2_{#31}    & #2_{#32}    & \cdots & #2_{#3#4}
  \end{pmatrix}
  \]
}

% 简化版命令（默认矩阵名为A，元素符号为a）: \quickmatrix{行数}{列数}
\newcommand{\quickmatrix}[2]{\flexmatrix{A}{a}{#1}{#2}}

\begin{document}
\title{习题 8.1}
\author{张志聪}
\maketitle

\section*{4}

设任意函数$f(x)$在区间$I$中含有第一类间断点。

反证法，假设$f$有原函数$F$，于是由原函数的定义可知，
$F$在区间$I$上是可导的，
所以$F$的导函数不会有第一类间断点（6-1-comment.tex中有详细证明），
这与$f$有第一类间断点矛盾。

\section*{5}

\begin{itemize}
  \item (12)

        \begin{align*}
          \int \sqrt{x\sqrt{x\sqrt{x}}} dx
           & = \int x^{\frac{1}{2}} x^{\frac{1}{4}} x^{\frac{1}{8}} dx \\
           & = \int x^{\frac{7}{8}}  dx                                \\
           & = \frac{8}{15} x^{\frac{15}{8}} + C
        \end{align*}


  \item (15)

        积化和差
        \begin{align*}
          \int cos x \cdot cos 2x dx
           & = \int \frac{1}{2} (cos 3 x + cos x) dx         \\
           & = \frac{1}{2} \int (cos 3 x + cos x) dx         \\
           & = \frac{1}{2} (\frac{1}{3} sin 3 x + sin x) + C
        \end{align*}

        另一种解法
        \begin{align*}
          \int cos x \cdot cos 2x dx
           & = \int cos x (cos^2x - sin^2 x) dx \\
           & = \int cos x (1 - 2 sin^2 x) dx    \\
           & = \int cos x - 2 cos x sin^2 x dx  \\
           & = sinx - \frac{2}{3} sin^3 x + C
        \end{align*}

  \item (18)

        可知$x \neq 0$。
        \begin{align*}
          \int \frac{\sqrt{x^4 + x^{-4} + 2}}{x^3} dx
           & = \int \frac{\sqrt{(x^2 + x^{-2})^2}}{x^3} dx \\
           & = \int \frac{x^2 + x^{-2}}{x^3} dx            \\
           & = \int \frac{1}{x} + \frac{1}{x^5} dx         \\
           & = \int \frac{1}{x} dx + \int \frac{1}{x^5} dx \\
           & = \ln |x| - \frac{1}{4} x^{-4} + C
        \end{align*}

\end{itemize}

\section*{6}

仿照例6

\begin{itemize}
  \item (2) $\int |sinx| dx$.

        设
        \begin{equation*}
          f(x) =
          \begin{cases}
            sinx,  & x \in [2k \pi, \pi + 2k \pi]        \\
            -sinx, & x \in [\pi + 2k \pi, 2\pi + 2k \pi]
          \end{cases}
          , k \in \mathbb{Z}
        \end{equation*}
        因为$f(x) = \sqrt{sin^2x}$是初等函数，所以在$\mathbb{R}$上连续，
        于是不定积分$\int f(x) dx$在$\mathbb{R}$上存在。

        (i)$k = 0$时，$cosx$是$f(x)$在$[\pi, 2\pi]$上的一个原函数，
        设$F(x)$是$f(x)$的一个原函数，且满足
        \begin{align*}
          F(x) = cos x, x \in [\pi, 2\pi]
        \end{align*}
        则当$x \in [0, \pi]$时，$F'(x) = sinx$。所以存在常数$C_1$，使得
        \begin{align*}
          F(x) = -cosx + C_1
        \end{align*}
        因为$F(x)$在$\mathbb{R}$上可导，所以在$x = \pi$处连续，从而有
        \begin{align*}
          -1 & = 1 + C_1
        \end{align*}
        即$C_1 = -2$，因此
        \begin{equation*}
          F(x) =
          \begin{cases}
            -cosx - 2, & x \in [0, \pi]    \\
            cosx,      & x \in [\pi, 2\pi]
          \end{cases}
        \end{equation*}
        (ii)$k = 1$时，当$x \in [2\pi, 3\pi]$时$F'(x) = sin x$。所以存在常数$C_2$，使得
        \begin{align*}
          F(x) = -cos x + C_1, x \in [1\pi, 3\pi]
        \end{align*}
        因为$F(x)$在$\mathbb{R}$上可导，所以在$x = 2\pi$处连续，从而有
        \begin{align*}
          cos 2\pi = -cos 2\pi + C_2 \\
          1 = -1 + C_2
        \end{align*}
        即$C_2 = 2$。

        同理，当$x \in [3\pi, 4\pi]$时$F'(x) = -sin x$。所以存在常数$C_2$，使得
        \begin{align*}
          F(x) = cosx + C_3, x \in [3\pi, 4\pi]
        \end{align*}
        因为$F(x)$在$\mathbb{R}$上可导，所以在$x = 3\pi$处连续，从而有
        \begin{align*}
          -cos 3\pi + 2 = cos 3\pi + C_3 \\
          1 + 2 = -1 + C_3
        \end{align*}
        即$C_3 = 4$。
        因此
        \begin{equation*}
          F(x) =
          \begin{cases}
            -cosx + 2, & x \in [2\pi, 3\pi] \\
            cosx + 4,  & x \in [3\pi, 4\pi]
          \end{cases}
        \end{equation*}

        于是，$k$是非负整数时，我们假设
        \begin{equation*}
          F(x) =
          \begin{cases}
            -cosx - 2 + 4k, & x \in [2k \pi, \pi + 2k \pi]        \\
            cosx + 4k,      & x \in [\pi + 2k \pi, 2\pi + 2k \pi]
          \end{cases}
          k \in \mathbb{N}
        \end{equation*}
        对$k$进行归纳。\\
        已知$k = 0$时（其实$k=1$时，通过上面的证明可知也是成立的），命题成立。\\
        归纳假设$k-1$时，命题成立，即
        \begin{equation*}
          F(x) =
          \begin{cases}
            -cosx - 2 + 4(k-1), & x \in [2(k-1) \pi, \pi + 2(k-1) \pi]        \\
            cosx + 4(k-1),      & x \in [\pi + 2(k-1) \pi, 2\pi + 2(k-1) \pi]
          \end{cases}
        \end{equation*}
        成立。

        $k$时，
        当$x \in [2k \pi, \pi + 2k \pi]$时$F'(x) = sin x$。所以存在常数$C_k$，使得
        \begin{align*}
          F(x) = -cos x + C_k, x \in [2k \pi, \pi + 2k \pi]
        \end{align*}
        因为$F(x)$在$\mathbb{R}$上可导，所以在$x = 2k\pi$处连续，从而有
        \begin{align*}
          cos 2k\pi + 4(k-1) = -cos 2k\pi + C_k \\
          1 + 4k - 4 = -1 + C_k
        \end{align*}
        即$C_k = 4k -2$。

        同理，当$x \in [\pi + 2k \pi, 2\pi + 2k \pi]$时$F'(x) = -sin x$。所以存在常数$C_k$，使得
        \begin{align*}
          F(x) = cosx + C_k, x \in [\pi + 2k \pi, 2\pi + 2k \pi]
        \end{align*}
        因为$F(x)$在$\mathbb{R}$上可导，所以在$x = \pi + 2k \pi$处连续，从而有
        \begin{align*}
          -cos (\pi + 2k \pi) + 4k -2 = cos(\pi + 2k \pi) + C_k \\
          1 + 4k - 2 = -1 + C_k
        \end{align*}
        即$C_k = 4k$。
        因此
        \begin{equation*}
          F(x) =
          \begin{cases}
            -cosx - 2 + 4k, & x \in [2k \pi, \pi + 2k \pi]        \\
            cosx + 4k,      & x \in [\pi + 2k \pi, 2\pi + 2k \pi]
          \end{cases}
        \end{equation*}
        归纳完成。

        于是，$k$是非正整数时，令$n = -k$，我们假设
        \begin{equation*}
          F(x) =
          \begin{cases}
            -cosx - 2 - 4n, & x \in [-2n \pi, \pi - 2n \pi]       \\
            cosx - 4n,      & x \in [\pi - 2n \pi, 2\pi - 2n \pi]
          \end{cases}
        \end{equation*}
        对$n$进行归纳假设。

        $n = 0$时，命题成立。

        归纳假设$n-1$时，命题成立，即
        \begin{equation*}
          F(x) =
          \begin{cases}
            -cosx - 2 - 4(n-1), & x \in [-2(n-1) \pi, \pi - 2(n-1) \pi]       \\
            cosx - 4(n-1),      & x \in [\pi - 2(n-1) \pi, 2\pi - 2(n-1) \pi]
          \end{cases}
        \end{equation*}

        $n$时，
        当$x \in [\pi - 2n \pi, 2\pi - 2n \pi]$时$F'(x) = -sin x$。所以存在常数$C_n$，使得
        \begin{align*}
          F(x) = cos x + C_n, x \in [\pi - 2n \pi, 2\pi - 2n \pi]
        \end{align*}
        因为$F(x)$在$\mathbb{R}$上可导，所以在$x = 2\pi - 2n \pi$处连续，从而有
        \begin{align*}
          -cos(2\pi - 2n \pi) - 2 - 4(n-1) = cos (2\pi - 2n \pi) + C_n \\
          -1 - 2 - 4n + 4 = 1 + C_n
        \end{align*}
        即$C_n = - 4n$。

        同理，当$x \in [-2n \pi, \pi - 2n \pi]$时$F'(x) = sin x$。所以存在常数$C_n$，使得
        \begin{align*}
          F(x) = -cos x + C_n, x \in [-2n \pi, \pi - 2n \pi]
        \end{align*}
        因为$F(x)$在$\mathbb{R}$上可导，所以在$x = \pi - 2n \pi$处连续，从而有
        \begin{align*}
          cos (\pi - 2n \pi) - 4n = -cos (\pi - 2n \pi) + C_n \\
          -1 + 4n = 1 + C_n
        \end{align*}
        即$C_n = 4n - 2$。
        因此
        \begin{equation*}
          F(x) =
          \begin{cases}
            -cosx - 2 - 4n, & x \in [-2n \pi, \pi - 2n \pi]       \\
            cosx - 4n,      & x \in [\pi - 2n \pi, 2\pi - 2n \pi]
          \end{cases}
        \end{equation*}
        归纳完成。

        令$k = -n$，于是，$k$为负整数时，我们有
        \begin{equation*}
          F(x) =
          \begin{cases}
            -cosx - 2 + 4k, & x \in [2k \pi, \pi + 2k \pi]       \\
            cosx + 4k,      & x \in [\pi + 2k \pi, 2\pi + 2k \pi]
          \end{cases}
        \end{equation*}

        综上，
        \begin{equation*}
          F(x) =
          \begin{cases}
            -cosx - 2 + 4k, & x \in [2k \pi, \pi + 2k \pi]        \\
            cosx + 4k,      & x \in [\pi + 2k \pi, 2\pi + 2k \pi]
          \end{cases}
          k \in \mathbb{Z}
        \end{equation*}
        

\end{itemize}



\section*{7}

\begin{zremark}
  P94页的注2的理解对这道题很有帮助。

  复合函数$f'(u), u = \varphi (x)$其实是对$u$求导，而$(f(u))'$是对$x$求导，
  两者的自变量不同，前者的自变量是$u$，后者是$x$。一个更好理解的方式是：
  前者就是对原函数$f$的求导，而后者是对复合函数$f \circ \varphi$的求导。
\end{zremark}

令$y = arctan x, x = tan y$，于是
\begin{align*}
  f'(arctan x) & = x^2     \\
  f'(y)        & = tan^2 y
\end{align*}
因为函数与自变量的符号无关，所以
\begin{align*}
  f'(x) = tan^2 x
\end{align*}
我们有
\begin{align*}
  f(x) & = \int tan^2 x dx \\
       & = tanx - x + C
\end{align*}

\section*{8}

 (i)$f$有原函数$F$，但$f$有第二类间断点的情况
\begin{equation*}
  F(x) =
  \begin{cases*}
    x^2 sin \frac{1}{x}, x \neq 0 \\
    0, x = 0
  \end{cases*}
\end{equation*}

证明略

(ii) $f$没有原函数，但有第二类间断点的情况
\begin{equation*}
  f(x) =
  \begin{cases*}
    \frac{1}{x}, x \neq 0 \\
    1, x = 0
  \end{cases*}
\end{equation*}

\textbf{证明}

假设$f$有原函数$F$，于是，我们有
\begin{align*}
  F'(1) = f(1) = 1 \\
  F'(-1) = f(-1) = -1
\end{align*}
由定理6.5（达布定理）可知，存在$\xi \in (-1, 1)$，使得
\begin{align*}
  F'(\xi) = 0
\end{align*}
即
\begin{align*}
  f(\xi) = 0
\end{align*}
这与$f(x) \neq 0$矛盾。



\end{document}